% !PRINTED 558
\weberpage{558}
%<<CONT
tative 1. To define opposite classes, we take a representative $\varphi$ of a class $A$ and a number $\alpha$ divisible by $\varphi$. If now $\alpha = \varphi\chi$, then $\chi$ is a representative of the class $A^{-1}$, and we have $AA^{-1} = O$.

If $A$ is an arbitrary class, and $h$ the class number, then always

\[ A^h = O, \]

and if $k$ is the least positive number that satisfies the condition

\[ A^k = O \]

then $k$ is a divisor of $h$. From this results the theorem:

3.~\gesperrt{Every functional $\varphi$ in $\Omega$ belongs to a definite exponent $k$, which is a divisor of the class number $h$, so that $\varphi^k$ is associated with a number in $\Omega$.}

\begin{center}
\rule{4cm}{0.4pt}
\end{center}
